\documentclass[12pt,a4paper]{article}
\usepackage[utf8]{inputenc}
\usepackage[T1]{fontenc}
\usepackage[english,romanian]{babel}
\usepackage{graphicx}
\usepackage{subcaption}
\usepackage{amsmath,amssymb}
\usepackage{fancyhdr}
\usepackage[a4paper,margin=2.5cm]{geometry}
\usepackage{dirtytalk}
\usepackage{csquotes}
\usepackage[hidelinks]{hyperref}
\usepackage{listings}
\usepackage{xcolor}
\usepackage{array}
\usepackage{booktabs}
\usepackage{tabularx}

\lstset{
    language=Python,
    basicstyle=\ttfamily\small,
    numbers=left,
    numberstyle=\tiny,
    frame=single,
    breaklines=true,
    showstringspaces=false,
    captionpos=b
}
% listings nu înțelege direct diacriticele UTF-8, de aceea le traducem explicit.
\lstset{
    literate=
        {Ă}{{\u{A}}}1 {ă}{{\u{a}}}1
        {Â}{{\^A}}1 {â}{{\^a}}1
        {Î}{{\^I}}1 {î}{{\^i}}1
        {Ș}{{\textcommabelow{S}}}1 {ș}{{\textcommabelow{s}}}1
        {Ț}{{\textcommabelow{T}}}1 {ț}{{\textcommabelow{t}}}1
}
\lstdefinestyle{terminal}{
    language={},
    numbers=none,
    backgroundcolor=\color{black!5},
    basicstyle=\ttfamily\scriptsize
}
\lstdefinestyle{config}{
    language={},
    numbers=none,
    basicstyle=\ttfamily\footnotesize
}
\setlength{\parindent}{1.25cm}
\setlength{\parskip}{0.25em}

\begin{document}
\begin{titlepage}
\begin{center}
    % Înlocuiți calea dacă logoul este stocat în altă parte.
    \includegraphics[width=0.3\textwidth]{lab1/image .png}\\[0.2cm]
    \textbf{MINISTERUL EDUCAȚIEI, CULTURII ȘI CERCETĂRII\\
    AL REPUBLICII MOLDOVA}\\[0.3cm]
    \textbf{Universitatea Tehnică a Moldovei\\
    Facultatea Calculatoare, Informatică și Microelectronică\\
    Departamentul Inginerie Software și Automatică}\\[2cm]
    \Huge\textbf{Raport}\\[0.5cm]
    \Large Lucrarea de laborator nr. 5\\[0.3cm]
    \large\textbf{la Criptografie și Securitate}\\[0.8cm]
    \textbf{Ciur Alexandrina, FAF-241}\\[2.5cm]
    \begin{flushright}
        \textit{Verificat de:}\\
        \textbf{Zaica M.}\\
        DISA, FCIM, UTM
    \end{flushright}
    \vfill
    Chișinău -- 2026
\end{center}
\end{titlepage}

\section*{Tema: Infrastructura cu chei publice (PKI) și semnătura digitală}

\section{Scopul lucrării}

Studierea infrastructurii cu chei publice (PKI) și a semnăturii digitale și
crearea, cu ajutorul utilitarului \textbf{OpenSSL}, a unei PKI interne: o
autoritate de certificare (CA) cu certificat autosemnat, care emite și revocă
certificatele utilizatorilor, iar utilizatorii își semnează fișierele și
verifică semnăturile.

\section{Sarcina lucrării}

\begin{quote}
\itshape
Să se creeze o PKI internă folosind utilitarul OpenSSL. Este necesară generarea
cheii private root și inițializarea unei autorități de certificare (CA). Pentru
CA se creează un certificat autosemnat. Sistemul trebuie să poată emite și
revoca cheia privată a utilizatorului, pentru ca acesta să poată genera ulterior
o semnătură digitală. Fiecare utilizator care obține o semnătură poate semna un
document/fișier și poate verifica această semnătură.
\end{quote}

Cerințe:
\begin{itemize}
    \item pentru generarea cheilor private se folosește algoritmul \textbf{RSA};
    \item cheile utilizatorilor au valabilitatea de \textbf{365 de zile} și
    dimensiunea de cel puțin \textbf{2048 de biți};
    \item cheia privată a CA are \textbf{4096 de biți}, iar certificatul
    autosemnat expiră după \textbf{10 ani (3650 de zile)}.
\end{itemize}

\section{Considerații teoretice}

\subsection{Criptografia cu chei publice}

În criptografia asimetrică fiecare participant are o pereche de chei: o
\textbf{cheie privată}, pe care o păstrează doar el, și o \textbf{cheie publică},
pe care o poate face cunoscută tuturor. Ce s-a făcut cu una dintre chei poate fi
verificat sau desfăcut doar cu cealaltă. Pentru semnătura digitală, semnatarul
folosește cheia sa privată, iar oricine poate verifica semnătura cu cheia
publică.

Rămâne însă o problemă: cum știe verificatorul că o anumită cheie publică
aparține într-adevăr persoanei respective și nu unui atacator care s-a dat drept
ea? Această problemă este rezolvată de infrastructura cu chei publice.

\subsection{Infrastructura cu chei publice (PKI)}

\textbf{PKI} (\textit{Public Key Infrastructure}) este ansamblul de componente,
reguli și proceduri prin care se leagă o cheie publică de identitatea
proprietarului ei. Componentele principale sunt:
\begin{itemize}
    \item \textbf{autoritatea de certificare (CA)} -- entitatea de încredere care
    verifică identitatea solicitanților și semnează certificatele lor cu cheia
    sa privată;
    \item \textbf{certificatul digital} -- un document electronic (standardul
    \textbf{X.509}) care conține cheia publică a titularului, numele lui
    (\textit{subject}), numele CA emitente (\textit{issuer}), numărul de serie,
    perioada de valabilitate, extensiile și semnătura CA;
    \item \textbf{cererea de certificat (CSR)} -- mesajul prin care utilizatorul
    își trimite cheia publică și numele spre semnare la CA. CSR-ul este semnat
    cu cheia privată a utilizatorului, ceea ce dovedește că el deține perechea
    cheii publice, dar cheia privată nu părăsește niciodată utilizatorul;
    \item \textbf{lista certificatelor revocate (CRL)} -- o listă semnată de CA
    care conține numerele de serie ale certificatelor anulate înainte de
    expirare (de exemplu, din cauza compromiterii cheii private);
    \item \textbf{baza de date a CA} -- evidența tuturor certificatelor emise și
    a stării lor (valid, revocat, expirat).
\end{itemize}

\subsection{Certificatul autosemnat și lanțul de încredere}

Certificatul CA root este \textbf{autosemnat}: câmpurile \textit{subject} și
\textit{issuer} coincid, iar el este semnat cu propria cheie privată a CA. El
nu poate fi verificat de nimeni de deasupra, de aceea este distribuit în
prealabil tuturor participanților și considerat de încredere. Certificatul unui
utilizator este semnat de CA, deci poate fi verificat cu cheia publică din
certificatul CA. Astfel se formează \textbf{lanțul de încredere}: CA root
$\rightarrow$ utilizator. Extensia \texttt{basicConstraints = CA:TRUE} marchează
certificatul care are dreptul să semneze alte certificate, iar
\texttt{keyUsage} limitează scopurile în care poate fi folosită cheia.

Deoarece toate certificatele depind de cheia privată root, aceasta trebuie
protejată cel mai bine: are o dimensiune mai mare (4096 de biți) și este
păstrată criptată cu o parolă.

\subsection{Semnătura digitală}

Semnătura digitală asigură:
\begin{itemize}
    \item \textbf{integritatea} -- orice modificare a documentului face semnătura
    invalidă;
    \item \textbf{autentificarea} -- se cunoaște cine a semnat documentul;
    \item \textbf{nerepudierea} -- semnatarul nu poate nega ulterior că a semnat,
    deoarece doar el deține cheia privată.
\end{itemize}

Documentul nu se semnează direct, ci se semnează \textbf{amprenta} (rezumatul)
lui, calculată cu o funcție hash. Pentru semnătura RSA cu SHA-256:
\begin{enumerate}
    \item \textbf{semnarea}: se calculează $h = \mathrm{SHA256}(M)$, $h$ este
    completat după schema PKCS\#1 v1.5 până la un număr $m < n$, iar semnătura
    este $s = m^{d} \bmod n$, unde $(d, n)$ este cheia privată;
    \item \textbf{verificarea}: se calculează $m' = s^{e} \bmod n$ cu cheia
    publică $(e, n)$ și se compară hash-ul extras din $m'$ cu
    $\mathrm{SHA256}(M)$ calculat din nou de verificator.
\end{enumerate}
Dimensiunea semnăturii este egală cu dimensiunea modulului $n$: pentru o cheie
de 2048 de biți semnătura are 256 de octeți.

\subsection{Algoritmul DSA și algoritmul RSA}

\textbf{DSA} (\textit{Digital Signature Algorithm}) este algoritmul de semnătură
standardizat în SUA prin FIPS 186. El se bazează pe dificultatea problemei
logaritmului discret, folosește parametrii publici $(p, q, g)$ și o valoare
aleatoare $k$ nouă la fiecare semnare; semnătura este perechea $(r, s)$. Spre
deosebire de RSA, DSA poate fi folosit doar pentru semnare, nu și pentru
criptare. Versiunea FIPS 186-5 (2023) nu mai permite generarea de semnături noi
cu DSA și recomandă RSA, ECDSA sau EdDSA.

Conform cerințelor lucrării, cheile au fost generate cu \textbf{RSA}, deci
semnătura digitală este RSA (\texttt{sha256WithRSAEncryption}). Securitatea RSA
se bazează pe dificultatea factorizării modulului $n = p \cdot q$.

\subsection{Revocarea certificatelor}

Un certificat poate deveni nesigur înainte de expirare, de exemplu dacă cheia
privată a fost furată sau dacă angajatul a părăsit organizația. În acest caz CA
\textbf{revocă} certificatul: îl marchează în baza de date ca revocat și
publică o nouă CRL care conține numărul lui de serie și motivul revocării
(RFC 5280). La verificarea unei semnături, verificatorul trebuie să controleze
nu doar semnătura, ci și dacă certificatul semnatarului este emis de CA, nu a
expirat și nu este în CRL.

Revocarea nu poate \enquote{șterge} cheia privată din calculatorul
utilizatorului, iar matematic aceasta poate semna în continuare. Însă orice
verificator care consultă CRL află că acea cheie nu mai este de încredere, deci
semnăturile ei nu mai sunt acceptate.

\subsection{Utilitarul OpenSSL}

\textbf{OpenSSL} este o bibliotecă și un utilitar în linie de comandă pentru
criptografie. În lucrare au fost folosite comenzile:

\begin{center}
\renewcommand{\arraystretch}{1.2}
\begin{tabularx}{\textwidth}{@{}lX@{}}
\toprule
\textbf{Comanda} & \textbf{Rolul} \\
\midrule
\texttt{genpkey} & generarea cheilor private RSA \\
\texttt{req} & crearea CSR-ului sau a certificatului autosemnat (\texttt{-x509}) \\
\texttt{ca} & semnarea CSR-urilor, revocarea certificatelor, generarea CRL \\
\texttt{verify} & verificarea certificatului față de CA și CRL \\
\texttt{x509} & afișarea certificatului, extragerea cheii publice \\
\texttt{dgst} & calcularea hash-ului, semnarea și verificarea semnăturii \\
\texttt{crl}, \texttt{pkey} & afișarea CRL, verificarea parolei cheii CA \\
\bottomrule
\end{tabularx}
\end{center}

\section{Implementarea}

Programul a fost realizat în limbajul \textbf{Python 3} (fișierul
\texttt{lab5/pki.py}), care apelează utilitarul \texttt{openssl} prin modulul
\texttt{subprocess}. Programul are un meniu și afișează fiecare comandă OpenSSL
executată, astfel încât se vede ce face PKI la fiecare pas. Toate fișierele PKI
se păstrează în folderul \texttt{lab5/pki/}:

\begin{lstlisting}[style=config, caption={Structura folderului PKI}]
pki/
  openssl.cnf          configuratia CA
  private/ca.key       cheia privata root (RSA 4096, criptata AES-256)
  certs/ca.crt         certificatul autosemnat al CA (3650 de zile)
  crl/ca.crl           lista certificatelor revocate
  index.txt            baza de date a CA (starea certificatelor)
  serial, crlnumber    urmatorul numar de serie / numar CRL
  newcerts/1000.pem    copiile certificatelor emise
  users/<nume>/        <nume>.key, <nume>.csr, <nume>.crt, <nume>.pub
\end{lstlisting}

Folderul \texttt{pki/} conține chei private, de aceea a fost adăugat în
\texttt{.gitignore} și nu este publicat în repozitoriu.

\subsection{Parametrii și configurația CA}

\begin{lstlisting}[caption={Parametrii impuși de sarcină}]
CA_KEY_BITS = 4096
CA_DAYS = 3650
USER_KEY_BITS = (2048, 3072, 4096)
USER_DAYS = 365
CRL_DAYS = 30
DIGEST = "sha256"
\end{lstlisting}

Utilizatorul poate alege doar dimensiuni de cheie de cel puțin 2048 de biți.
Programul generează fișierul \texttt{openssl.cnf}; secțiunile principale sunt:

\begin{lstlisting}[style=config, caption={Fragment din \texttt{openssl.cnf}}]
[ CA_default ]
database         = $dir/index.txt
private_key      = $dir/private/ca.key
certificate      = $dir/certs/ca.crt
default_md       = sha256
default_days     = 365
policy           = policy_user
unique_subject   = no

[ v3_ca ]
basicConstraints = critical, CA:true
keyUsage         = critical, keyCertSign, cRLSign, digitalSignature

[ usr_cert ]
basicConstraints = critical, CA:false
keyUsage         = critical, digitalSignature, nonRepudiation
\end{lstlisting}

Extensia \texttt{v3\_ca} se aplică certificatului CA, iar \texttt{usr\_cert}
certificatelor utilizatorilor: acestea nu pot semna alte certificate
(\texttt{CA:false}) și pot fi folosite doar pentru semnătură digitală.
Opțiunea \texttt{unique\_subject = no} permite emiterea unui certificat nou cu
același nume după revocarea celui vechi.

\subsection{Apelarea OpenSSL și protecția parolei CA}

\begin{lstlisting}[caption={Funcția care execută comenzile OpenSSL}]
def run_openssl(*args, password=None, quiet=False, check=True):
    command = [OPENSSL, *args]
    if not quiet:
        print("  $ openssl " + " ".join(_quote(arg) for arg in args))
    PKI_DIR.mkdir(exist_ok=True)
    environment = dict(os.environ)
    if password is not None:
        environment["CA_PASS"] = password
    result = subprocess.run(
        command, cwd=PKI_DIR, env=environment,
        stdin=subprocess.DEVNULL, capture_output=True, text=True,
    )
    if not check:
        return result
    if result.returncode != 0:
        message = (result.stderr or result.stdout).strip()
        if "bad decrypt" in message or "maybe wrong password" in message:
            raise PkiError("Parola cheii CA este gresita.")
        raise PkiError(message or f"openssl a terminat cu codul {result.returncode}")
    return result.stdout
\end{lstlisting}

Parola cheii CA se citește cu \texttt{getpass} (nu apare pe ecran) și se
transmite către OpenSSL prin variabila de mediu \texttt{CA\_PASS}
(\texttt{-passin env:CA\_PASS}), nu prin linia de comandă, unde ar putea fi
văzută de alți utilizatori în lista proceselor.

\subsection{Inițializarea CA}

\begin{lstlisting}[caption={Generarea cheii root și a certificatului autosemnat}]
run_openssl("genpkey", "-algorithm", "RSA",
            "-pkeyopt", f"rsa_keygen_bits:{CA_KEY_BITS}",
            "-aes-256-cbc", "-pass", "env:CA_PASS",
            "-out", CA_KEY, password=password)
_restrict(PKI_DIR / CA_KEY)

run_openssl("req", "-config", CONFIG, "-new", "-x509",
            "-key", CA_KEY, "-passin", "env:CA_PASS",
            "-days", str(CA_DAYS), f"-{DIGEST}",
            "-extensions", "v3_ca", "-subj", CA_SUBJECT,
            "-out", CA_CERT, password=password)

generate_crl(password)
\end{lstlisting}

Cheia root de 4096 de biți se salvează criptată cu AES-256, iar
\texttt{\_restrict} îi lasă drept de citire doar proprietarului. Opțiunea
\texttt{-x509} a comenzii \texttt{req} creează direct un certificat autosemnat
(în loc de o cerere), valabil 3650 de zile. La final se publică o CRL goală.

\subsection{Emiterea cheii și a certificatului pentru un utilizator}

\begin{lstlisting}[caption={Funcția \texttt{issue\_user}}]
def issue_user(name, bits, password):
    key, csr, certificate = user_paths(name)
    ensure_no_valid_certificate(name)
    check_ca_password(password)
    (PKI_DIR / "users" / name).mkdir(parents=True, exist_ok=True)

    run_openssl("genpkey", "-algorithm", "RSA",
                "-pkeyopt", f"rsa_keygen_bits:{bits}", "-out", key)
    _restrict(PKI_DIR / key)

    run_openssl("req", "-config", CONFIG, "-new", "-key", key,
                "-subj", USER_SUBJECT.format(name=name), "-out", csr)

    run_openssl("ca", "-config", CONFIG, "-batch", "-notext",
                "-extensions", "usr_cert", "-days", str(USER_DAYS),
                "-passin", "env:CA_PASS", "-in", csr, "-out", certificate,
                password=password)
    return certificate
\end{lstlisting}

Emiterea are trei pași: utilizatorul își generează cheia RSA, creează cererea
CSR, iar CA o semnează și emite certificatul pe 365 de zile. Comanda
\texttt{openssl ca} atribuie certificatului următorul număr de serie și îl
înregistrează în \texttt{index.txt}. Înainte de a genera ceva, programul
verifică parola CA (\texttt{check\_ca\_password}) și refuză emiterea dacă
utilizatorul are deja un certificat valid.

\subsection{Revocarea certificatului}

\begin{lstlisting}[caption={Funcția \texttt{revoke\_user}}]
def revoke_user(name, reason, password):
    _, _, certificate = user_paths(name)
    if not (PKI_DIR / certificate).exists():
        raise PkiError(f"Utilizatorul {name} nu are certificat.")
    check_ca_password(password)
    run_openssl("ca", "-config", CONFIG, "-revoke", certificate,
                "-crl_reason", reason, "-passin", "env:CA_PASS",
                password=password)
    generate_crl(password)
\end{lstlisting}

Utilizatorul alege motivul revocării dintr-o listă (\texttt{keyCompromise},
\texttt{affiliationChanged}, \texttt{superseded},
\texttt{cessationOfOperation}). Certificatul primește starea \texttt{R} în baza
de date, iar CA publică imediat o CRL nouă care îl conține.

\subsection{Verificarea stării certificatului}

\begin{lstlisting}[caption={Verificarea certificatului față de CA și CRL}]
def certificate_status(certificate):
    try:
        output = run_openssl("verify", "-CAfile", CA_CERT,
                             "-CRLfile", CA_CRL, "-crl_check",
                             str(certificate), quiet=True)
        return True, output.strip()
    except PkiError as error:
        return False, str(error)
\end{lstlisting}

\texttt{openssl verify} controlează semnătura CA de pe certificat, perioada de
valabilitate și, datorită opțiunii \texttt{-crl\_check}, dacă certificatul nu se
află în CRL. Această funcție este folosită atât la semnare, cât și la
verificare.

\subsection{Semnarea și verificarea fișierelor}

\begin{lstlisting}[caption={Semnarea unui fișier}]
def sign_file(name, path):
    key, _, certificate = user_paths(name)
    if not (PKI_DIR / key).exists():
        raise PkiError(f"Utilizatorul {name} nu are cheie privata.")
    valid, message = certificate_status(certificate)
    if not valid:
        raise PkiError(f"Certificatul lui {name} nu este valid, semnarea este "
                       f"interzisa.\n{message}")
    signature = path.with_name(f"{path.name}.{name}.sig")
    run_openssl("dgst", f"-{DIGEST}", "-sign", key,
                "-out", _arg_path(signature), _arg_path(path))
    return signature
\end{lstlisting}

\texttt{openssl dgst -sha256 -sign} calculează hash-ul SHA-256 al fișierului și
îl semnează cu cheia privată RSA. Semnătura se salvează lângă fișier, cu numele
\texttt{<fișier>.<utilizator>.sig}. Dacă certificatul utilizatorului a fost
revocat sau a expirat, programul nu permite semnarea.

\begin{lstlisting}[caption={Verificarea semnăturii (fragment)}]
certificate_ok, certificate_message = certificate_status(certificate)

public_key = f"users/{name}/{name}.pub"
run_openssl("x509", "-in", certificate, "-pubkey", "-noout", "-out", public_key)

result = run_openssl("dgst", f"-{DIGEST}", "-verify", public_key,
                     "-signature", _arg_path(signature), _arg_path(path),
                     check=False)
signature_ok = result.returncode == 0
\end{lstlisting}

Verificarea are trei pași: (1) se verifică certificatul semnatarului față de CA
și CRL; (2) cheia publică se extrage din certificat, nu se ia dintr-o sursă
oarecare, deci este garantată de CA; (3) \texttt{openssl dgst -verify}
verifică semnătura. Semnătura este acceptată doar dacă ambele verificări
reușesc. Dacă semnătura corespunde fișierului, dar certificatul este revocat,
programul o raportează separat ca semnătură în care nu se poate avea încredere.

\section{Rezultatele execuției}

Programul se lansează din folderul \texttt{lab5} (OpenSSL trebuie să fie
instalat și disponibil în \texttt{PATH}; altfel calea spre \texttt{openssl} se
poate indica prin variabila de mediu \texttt{OPENSSL}):
\begin{lstlisting}[style=terminal]
python pki.py
\end{lstlisting}

Pentru testare a fost folosit fișierul \texttt{document.txt}:
\begin{lstlisting}[style=terminal]
Contract de prestari servicii nr. 17
Suma: 5000 MDL
\end{lstlisting}

În listingurile de mai jos meniul principal a fost afișat doar o dată.

\subsection{Inițializarea CA}

\begin{lstlisting}[style=terminal, caption={Crearea CA și a certificatului autosemnat}]
=== PKI interna si semnatura digitala ===
1. Initializarea CA (cheia root + certificatul autosemnat)
2. Emiterea cheii si a certificatului pentru un utilizator
3. Revocarea certificatului unui utilizator
4. Semnarea unui fisier
5. Verificarea unei semnaturi
6. Afisarea certificatelor si a listei CRL
0. Iesire
Alege (0-6): 1
Parola noua pentru cheia CA (min. 4 caractere):
Repetati parola:

[1/3] Cheia privata root (RSA 4096 biti, criptata AES-256):
  $ openssl genpkey -algorithm RSA -pkeyopt rsa_keygen_bits:4096 -aes-256-cbc -pass env:CA_PASS -out private/ca.key

[2/3] Certificatul autosemnat al CA (3650 de zile):
  $ openssl req -config openssl.cnf -new -x509 -key private/ca.key -passin env:CA_PASS -days 3650 -sha256 -extensions v3_ca -subj "/C=MD/L=Chisinau/O=UTM/OU=FCIM FAF-241/CN=FAF-241 Root CA" -out certs/ca.crt

[3/3] Lista certificatelor revocate (CRL), deocamdata goala:
  $ openssl ca -config openssl.cnf -gencrl -passin env:CA_PASS -out crl/ca.crl

CA a fost creata. Certificatul autosemnat:
  subject=C=MD, L=Chisinau, O=UTM, OU=FCIM FAF-241, CN=FAF-241 Root CA
  issuer=C=MD, L=Chisinau, O=UTM, OU=FCIM FAF-241, CN=FAF-241 Root CA
  serial=1342055DB5EC9D13A3D6E7751CF38096E8FB5ABC
  notBefore=Nov 22 11:20:55 2026 GMT
  notAfter=Nov 19 11:20:55 2036 GMT
  X509v3 Basic Constraints: critical
      CA:TRUE
  X509v3 Key Usage: critical
      Digital Signature, Certificate Sign, CRL Sign
  Cheia publica: RSA 4096 biti
\end{lstlisting}

Câmpurile \textit{subject} și \textit{issuer} coincid, deci certificatul este
autosemnat. Cheia are 4096 de biți, iar între \texttt{notBefore} și
\texttt{notAfter} sunt 3650 de zile (10 ani).

\subsection{Emiterea certificatelor pentru utilizatori}

Au fost creați doi utilizatori: \texttt{alexandrina} (cheie de 2048 de biți) și
\texttt{ion} (cheie de 3072 de biți).

\begin{lstlisting}[style=terminal, caption={Emiterea cheii și a certificatului pentru un utilizator}]
Alege (0-6): 2
Numele utilizatorului: alexandrina
Dimensiunea cheii: 1. 2048 biti, 2. 3072 biti, 3. 4096 biti
Alege (1-3): 1
Parola cheii CA:

[1/3] Cheia privata a utilizatorului (RSA 2048 biti):
  $ openssl genpkey -algorithm RSA -pkeyopt rsa_keygen_bits:2048 -out users/alexandrina/alexandrina.key

[2/3] Cererea de certificat (CSR):
  $ openssl req -config openssl.cnf -new -key users/alexandrina/alexandrina.key -subj "/C=MD/L=Chisinau/O=UTM/OU=FCIM FAF-241/CN=alexandrina" -out users/alexandrina/alexandrina.csr

[3/3] CA semneaza cererea (365 de zile):
  $ openssl ca -config openssl.cnf -batch -notext -extensions usr_cert -days 365 -passin env:CA_PASS -in users/alexandrina/alexandrina.csr -out users/alexandrina/alexandrina.crt

Cheia privata:  pki/users/alexandrina/alexandrina.key
Certificatul:   pki/users/alexandrina/alexandrina.crt
  subject=C=MD, L=Chisinau, O=UTM, OU=FCIM FAF-241, CN=alexandrina
  issuer=C=MD, L=Chisinau, O=UTM, OU=FCIM FAF-241, CN=FAF-241 Root CA
  serial=1000
  notBefore=Nov 22 11:21:02 2026 GMT
  notAfter=Nov 22 11:21:02 2027 GMT
  X509v3 Basic Constraints: critical
      CA:FALSE
  X509v3 Key Usage: critical
      Digital Signature, Non Repudiation
  Cheia publica: RSA 2048 biti

Alege (0-6): 2
Numele utilizatorului: ion
Dimensiunea cheii: 1. 2048 biti, 2. 3072 biti, 3. 4096 biti
Alege (1-3): 2
...
  serial=1001
  notBefore=Nov 22 11:21:09 2026 GMT
  notAfter=Nov 22 11:21:09 2027 GMT
  ...
  Cheia publica: RSA 3072 biti
\end{lstlisting}

Certificatul utilizatorului este emis de CA (\textit{issuer} este
\texttt{FAF-241 Root CA}), este valabil exact 365 de zile și nu are dreptul să
semneze alte certificate (\texttt{CA:FALSE}).

\subsection{Semnarea și verificarea unui fișier}

\begin{lstlisting}[style=terminal, caption={Semnarea documentului și verificarea reușită}]
Alege (0-6): 4
Utilizatori: alexandrina, ion
Numele utilizatorului: alexandrina
Fisierul de semnat: document.txt
  $ openssl dgst -sha256 -sign users/alexandrina/alexandrina.key -out ../document.txt.alexandrina.sig ../document.txt

Semnatura (256 octeti) a fost salvata in: document.txt.alexandrina.sig

Alege (0-6): 5
Utilizatori: alexandrina, ion
Numele utilizatorului: alexandrina
Fisierul semnat: document.txt
Fisierul semnaturii [document.txt.alexandrina.sig]:

[1/3] Verificarea certificatului semnatarului (CA + CRL):
  $ openssl verify -CAfile certs/ca.crt -CRLfile crl/ca.crl -crl_check users/alexandrina/alexandrina.crt
  users/alexandrina/alexandrina.crt: OK

[2/3] Extragerea cheii publice din certificat:
  $ openssl x509 -in users/alexandrina/alexandrina.crt -pubkey -noout -out users/alexandrina/alexandrina.pub

[3/3] Verificarea semnaturii:
  $ openssl dgst -sha256 -verify users/alexandrina/alexandrina.pub -signature ../document.txt.alexandrina.sig ../document.txt
  Verified OK

REZULTAT: semnatura este VALIDA. Fisierul a fost semnat de alexandrina si nu a fost modificat.
\end{lstlisting}

Semnătura are 256 de octeți, adică 2048 de biți, cât modulul cheii RSA a
utilizatorului.

\subsection{Verificarea unui document modificat și a unui semnatar greșit}

Fișierul \texttt{document\_modificat.txt} este o copie a documentului în care
suma a fost schimbată din \texttt{5000} în \texttt{50000 MDL}. Se verifică
semnătura originală a documentului, mai întâi pentru fișierul modificat, apoi
ca și cum ar fi semnătura lui \texttt{ion}.

\begin{lstlisting}[style=terminal, caption={Semnătura nu corespunde fișierului modificat sau altui utilizator}]
Alege (0-6): 5
Numele utilizatorului: alexandrina
Fisierul semnat: document_modificat.txt
Fisierul semnaturii [document_modificat.txt.alexandrina.sig]: document.txt.alexandrina.sig
...
[1/3] Verificarea certificatului semnatarului (CA + CRL):
  users/alexandrina/alexandrina.crt: OK
...
[3/3] Verificarea semnaturii:
  $ openssl dgst -sha256 -verify users/alexandrina/alexandrina.pub -signature ../document.txt.alexandrina.sig ../document_modificat.txt
  Verification failure

REZULTAT: semnatura este INVALIDA (fisierul a fost modificat sau semnatura nu apartine acestui utilizator).

Alege (0-6): 5
Numele utilizatorului: ion
Fisierul semnat: document.txt
Fisierul semnaturii [document.txt.ion.sig]: document.txt.alexandrina.sig
...
  users/ion/ion.crt: OK
...
  $ openssl dgst -sha256 -verify users/ion/ion.pub -signature ../document.txt.alexandrina.sig ../document.txt
  Verification failure

REZULTAT: semnatura este INVALIDA (fisierul a fost modificat sau semnatura nu apartine acestui utilizator).
\end{lstlisting}

Schimbarea unei singure cifre modifică complet hash-ul SHA-256, de aceea
semnătura nu mai corespunde (\emph{integritatea}). Cheia publică a lui
\texttt{ion} nu poate verifica o semnătură făcută cu cheia privată a lui
\texttt{alexandrina} (\emph{autentificarea}).

\subsection{Revocarea certificatului}

Se presupune că cheia privată a utilizatorului \texttt{alexandrina} a fost
compromisă și certificatul ei este revocat.

\begin{lstlisting}[style=terminal, caption={Revocarea certificatului și efectul ei}]
Alege (0-6): 3
Utilizatori: alexandrina, ion
Numele utilizatorului: alexandrina
1. cheia privata a fost compromisa
2. utilizatorul a parasit organizatia
3. certificatul a fost inlocuit
4. certificatul nu mai este necesar
Motivul revocarii (1-4): 1
Parola cheii CA:

[1/2] Revocarea certificatului in baza de date a CA:
  $ openssl ca -config openssl.cnf -revoke users/alexandrina/alexandrina.crt -crl_reason keyCompromise -passin env:CA_PASS

[2/2] Publicarea noii liste CRL:
  $ openssl ca -config openssl.cnf -gencrl -passin env:CA_PASS -out crl/ca.crl

Certificatul lui alexandrina a fost revocat (keyCompromise).

Alege (0-6): 5
Numele utilizatorului: alexandrina
Fisierul semnat: document.txt
Fisierul semnaturii [document.txt.alexandrina.sig]:

[1/3] Verificarea certificatului semnatarului (CA + CRL):
  $ openssl verify -CAfile certs/ca.crt -CRLfile crl/ca.crl -crl_check users/alexandrina/alexandrina.crt
  C=MD, L=Chisinau, O=UTM, OU=FCIM FAF-241, CN=alexandrina
  error 23 at 0 depth lookup: certificate revoked
  error users/alexandrina/alexandrina.crt: verification failed
...
[3/3] Verificarea semnaturii:
  Verified OK

REZULTAT: semnatura corespunde fisierului, dar NU este de incredere: certificatul lui alexandrina nu este valid (revocat sau expirat).

Alege (0-6): 4
Numele utilizatorului: alexandrina
Fisierul de semnat: document.txt

EROARE: Certificatul lui alexandrina nu este valid, semnarea este interzisa.
C=MD, L=Chisinau, O=UTM, OU=FCIM FAF-241, CN=alexandrina
error 23 at 0 depth lookup: certificate revoked
error users/alexandrina/alexandrina.crt: verification failed
\end{lstlisting}

Matematic semnătura este în continuare corectă (\texttt{Verified OK}), însă
OpenSSL găsește certificatul în CRL (\texttt{error 23: certificate revoked}),
deci semnătura nu mai este acceptată, iar semnarea de noi fișiere este
blocată.

\subsection{Starea PKI}

\begin{lstlisting}[style=terminal, caption={Certificatele emise și lista CRL}]
Alege (0-6): 6
...
Certificatele emise (baza de date index.txt):
  1000  revocat  expira 2027-11-22 11:21  /C=MD/L=Chisinau/O=UTM/OU=FCIM FAF-241/CN=alexandrina  (revocat 2026-11-22 11:21, keyCompromise)
  1001  valid    expira 2027-11-22 11:21  /C=MD/L=Chisinau/O=UTM/OU=FCIM FAF-241/CN=ion

Lista CRL (crl/ca.crl):
  lastUpdate=Nov 22 11:21:40 2026 GMT
  nextUpdate=Dec 22 11:21:40 2026 GMT
  crlNumber=0x1001
  Certificate revocate: 1000
\end{lstlisting}

După revocare, utilizatorul poate primi o cheie și un certificat nou
(opțiunea~2), care primește următorul număr de serie (\texttt{1002}); cât timp
utilizatorul are un certificat valid, programul refuză emiterea unuia nou.

\subsection{Tratarea erorilor}

\begin{lstlisting}[style=terminal, caption={Parolă greșită, certificat existent, opțiune inexistentă}]
Alege (0-6): 3
Numele utilizatorului: ion
Motivul revocarii (1-4): 1
Parola cheii CA:

EROARE: Parola cheii CA este gresita.

Alege (0-6): 2
Numele utilizatorului: ion

EROARE: Utilizatorul ion are deja un certificat valid. Revocati-l inainte de a emite unul nou.

Alege (0-6): 9
Optiune invalida. Alegeti un numar intre 0 si 6.
\end{lstlisting}

Parola CA se verifică înainte de orice modificare, deci o parolă greșită nu
lasă în PKI fișiere create pe jumătate. Programul verifică și numele
utilizatorului (2--32 de caractere: litere latine, cifre, \texttt{\_},
\texttt{-}), iar operațiile 2--6 nu sunt permise înainte de inițializarea CA.

% Pentru a adăuga capturi de ecran, încărcați imaginile în Overleaf și
% decomentați blocul de mai jos (schimbați numele fișierului dacă e altul).
% \begin{figure}[h!]
%     \centering
%     \includegraphics[width=0.9\textwidth]{lab5/semnare.png}
%     \caption{Semnarea și verificarea unui fișier}
% \end{figure}

\section{Concluzii}

În cadrul acestei lucrări de laborator a fost studiată infrastructura cu chei
publice și semnătura digitală și a fost creată o PKI internă cu ajutorul
utilitarului OpenSSL, controlată de un program în Python. A fost generată cheia
root RSA de 4096 de biți, păstrată criptată cu AES-256, și certificatul
autosemnat al CA, valabil 3650 de zile. Programul emite pentru utilizatori chei
RSA de cel puțin 2048 de biți și certificate semnate de CA pentru 365 de zile,
revocă certificatele și publică lista CRL, semnează fișiere cu RSA și SHA-256 și
verifică semnăturile.

Experimentele au confirmat proprietățile semnăturii digitale: modificarea unei
singure cifre din document face semnătura invalidă, iar semnătura unui
utilizator nu poate fi verificată cu cheia publică a altuia. Am observat că
verificarea matematică a semnăturii nu este suficientă: după revocare, semnătura
veche rămâne corectă matematic, dar nu mai poate fi acceptată. De aceea o
verificare completă controlează și certificatul semnatarului: dacă este emis de
CA de încredere, dacă nu a expirat și dacă nu este în CRL.

Securitatea întregii PKI depinde de cheia privată a CA: cine o deține poate
emite certificate în numele oricui. În practică, cheia root se păstrează offline
sau într-un modul hardware (HSM), iar certificatele utilizatorilor sunt semnate
de o CA intermediară. Deși tema lucrării menționează algoritmul DSA, conform
cerințelor a fost folosit RSA, care este și recomandarea actuală a standardului
FIPS 186-5.

Codul sursă al lucrării este disponibil în repozitoriul GitHub:\\
\url{https://github.com/alecsyyyk/Cryptography-and-Security}

\end{document}
